\labsection{5}{Поиск экстремума функции одной переменной с помощью генетического алгоритма с бинарным представлением особей}

\labpart{Цель работы}
Изучить основные понятия теории генетических алгоритмов, реализовать классический простой ГА с
представлением решений в форме бинарных строк, классическими операторами кроссинговера и мутации и
исследовать влияние параметров ГА на точность и надёжность поиска.

\labpart{Теоретические сведения}
Раздел~\ref{sec:ga} и подразделы~\ref{sec:operators}.1--\ref{sec:operators}.4 теоретической части:
основные понятия ГА, простой ГА и пример его работы, двоичное кодирование вещественного числа
(\ref{eq:decode}) и выбор длины хромосомы (\ref{eq:chrom-length}), код Грея, пропорциональная селекция,
одноточечный кроссинговер и побитовая мутация, элитизм.

\labpart{Постановка задачи}
Дана функция одной переменной $y = f(t)$ (таблица~\ref{tab:l5-var}). Найти её экстремум на отрезке
$t \in [t_0; t_1]$ генетическим алгоритмом с бинарным представлением особей и графически
проиллюстрировать динамику поиска.

\begin{table}[!htb]
\caption{Варианты заданий к лабораторной работе №5}
\label{tab:l5-var}
\centering\small\setlength{\tabcolsep}{3.5pt}
\begin{tabular}{clc|clc}
\toprule
№ & Задача & $t$ & № & Задача & $t$ \\
\midrule
1 & $\max\ (1,5t + 0,9)\sin(\pi t + 1,1)$      & $[0; 5]$   & 5 & $\min\ (2,5t + 1,7)\sin(1,1\pi t + 0,7)$ & $[0; 10]$ \\
2 & $\min\ (1,3t + 1,9)\cos(1,1\pi t - 1,5)$    & $[-6; 6]$  & 6 & $\min\ (0,5t - 1,4)\cos(0,5\pi t + 1,1)$ & $[-9; 9]$ \\
3 & $\max\ (t + 1,3)\sin(0,5\pi t + 1)$         & $[0; 7]$   & 7 & $\max\ (1,7t + 1,5)\sin(0,7\pi t - 1,4)$ & $[0; 8]$ \\
4 & $\max\ (1,1t - 1,7)\cos(\pi t + 1,5)$       & $[-9; 9]$  & 8 & $\max\ (0,7t - 1,7)\cos(0,5\pi t + 1,5)$ & $[-5; 5]$ \\
\bottomrule
\end{tabular}
\end{table}

\labpart{Требования к программе}
Исходные данные программы: размер популяции, длина хромосомы (число битов), вероятности кроссинговера и
мутации, максимальное число поколений. В процессе поиска отображаются: график функции и положение на
нём особей текущей популяции; графики лучшего и среднего значений приспособленности по поколениям;
вещественное значение лучшей особи текущей популяции и значение функции для неё.

\labpart{Порядок выполнения работы}
\begin{steps}
\item Построить график функции и найти эталонный экстремум численно: перебор по мелкой сетке
\code{np.linspace} и уточнение \code{scipy.optimize.minimize\_scalar(method='bounded')}.
\item Выбрать длину хромосомы по требуемой точности (\ref{eq:chrom-length}), реализовать кодирование и
декодирование (\ref{eq:decode}).
\item Реализовать простой ГА на NumPy (подраздел~\ref{sec:ga-numpy}): случайную начальную популяцию, пропорциональную селекцию (для функций со
знакопеременными значениями и для поиска минимума~--- со сдвигом приспособленности), одноточечный
кроссинговер с вероятностью $P_c$, побитовую мутацию с вероятностью $P_m$, элитизм.
\item Выполнить основной запуск с визуализацией динамики популяции.
\item Оценить надёжность: выполнить не менее 30 независимых запусков и определить долю успешных
(найден экстремум с заданной точностью, например $|t - t^*| < 10^{-3}$) и медиану числа поколений до
успеха.
\item Исследовать влияние размера популяции, длины хромосомы, $P_c$, $P_m$ и элитизма, меняя по одному
параметру.
\end{steps}

\labpart{Пример выполнения}
Вариант~3: максимум $f(t) = (t + 1,3)\sin(0,5\pi t + 1)$ на отрезке $[0; 7]$. Функция имеет локальный
максимум $f(0,572) = 1,772$ и глобальный, найденный перебором по сетке из 700\,001 точки с уточнением
\code{scipy.optimize.minimize\_scalar}: $t^* = 4,433776$, $f(t^*) = 5,698757$. Параметры ГА: $N = 60$,
$L = 20$ (шаг сетки $\Delta = 6,7 \cdot 10^{-6}$), $P_c = 0,85$, $P_m = 0,01$, 200 поколений, элитизм;
успех~--- $|t - t^*| < 10^{-3}$. Найдено $t = 4,433777$ с погрешностью $1,6 \cdot 10^{-6}$ на уровне
точности 20-битного кода; окрестность $10^{-3}$ достигнута на 10-м поколении (660 вычислений функции).
К 10-му поколению популяция покидает локальный максимум и собирается у глобального
(рисунок~\ref{fig:l5-pop}); средняя приспособленность отстаёт от лучшей из-за мутаций
(рисунок~\ref{fig:l5-conv}).

\begin{figure}[!htb]
\centering
\includegraphics[width=\textwidth]{\figdir/lab5_population.png}
\caption{Популяция простого ГА на графике функции варианта~3: поколения 0, 10 и 200}
\label{fig:l5-pop}
\end{figure}

\begin{figure}[!htb]
\centering
\includegraphics[width=\textwidth]{\figdir/lab5_convergence.png}
\caption{Слева~--- лучшее и среднее значения функции в основном запуске (пунктир~--- $f^*$); справа~---
лучшая особь поколения с элитизмом и без него (среднее по 30 запускам, $P_m = 0,05$)}
\label{fig:l5-conv}
\end{figure}

Из 50 независимых запусков 47 успешны, медиана~--- 13 поколений до попадания. Все три неудачи
остановились в одной точке $t = 4,375$ ($f = 5,674$) с кодом \bits{10011111111111111111}: ближайшие
лучшие решения начинаются с \bits{10100\ldots}, и переход к ним требует одновременно изменить
несколько старших битов~--- это хэмминговый обрыв (подраздел~\ref{sec:gray}). С кодом Грея те же
запуски успешны все 50 из 50 (медиана 11 поколений), а при $N = 20$~--- 30 из 30 против 12 из 30 у
двоичного кода. Влияние параметров показано в таблице~\ref{tab:l5-params}.

\begin{table}[!htb]
\caption{Влияние параметров простого ГА: доля успешных запусков, худшее $f$ и медиана поколений до
успеха (30 запусков на значение, двоичный код)}
\label{tab:l5-params}
\centering\small
\begin{tabular}{lcccc|lcccc}
\toprule
 & Знач. & Успех, \% & Худшее $f$ & Покол. & & Знач. & Успех, \% & Худшее $f$ & Покол. \\
\midrule
$N$ & 10  & 20  & 5,6741 & 49   & $P_m$ & 0     & 47  & 5,6740 & 8,5 \\
    & 20  & 40  & 5,6741 & 20   &       & 0,001 & 73  & 5,6741 & 13 \\
    & 60  & 93  & 5,6741 & 14   &       & 0,01  & 93  & 5,6741 & 14 \\
    & 100 & 100 & 5,6988 & 14,5 &       & 0,05  & 100 & 5,6988 & 11 \\
\midrule
$L$ & 8   & 0   & 5,6975 & ---  &       & 0,2   & 100 & 5,6988 & 17,5 \\
    & 12  & 100 & 5,6988 & 14   & $P_c$ & 0     & 93  & 5,6741 & 35 \\
    & 20  & 93  & 5,6741 & 14   &       & 0,5   & 97  & 5,6741 & 15 \\
    & 24  & 97  & 5,6741 & 18   &       & 0,95  & 100 & 5,6988 & 17,5 \\
\bottomrule
\end{tabular}
\end{table}

Размер популяции~--- главный фактор надёжности. При $L = 8$ шаг сетки $0,027$ больше допуска, и
попадание в окрестность $10^{-3}$ невозможно в принципе. Без мутации популяция вырождается, и больше
половины запусков неудачны; мутация $P_m \ge 0,05$ помогает выйти из ловушки у хэммингового обрыва.
Кроссинговер ускоряет поиск более чем вдвое (35 и 14 поколений при $P_c = 0$ и $0,85$). Без
элитизма найденный максимум теряется из-за мутаций: лучшая особь последнего поколения в среднем
$5,6931$ против $5,6988$.

\labpart{Контрольные вопросы}
\begin{questions}
\item Поясните понятия <<эволюционные вычисления>> и <<эволюционный поиск>>. Перечислите основные
парадигмы эволюционных вычислений.
\item Что такое генетический алгоритм? Назовите его отличительные особенности.
\item Какие компоненты нужно задать для определения ГА? Опишите простой ГА.
\item Поясните понятия <<приспособленность особи>> и <<приспособленность популяции>>. Как искать
минимум функции и работать с отрицательными значениями?
\item Опишите классические операторы репродукции, кроссинговера и мутации. Приведите примеры.
\item Выполните вручную одно поколение простого ГА для поиска максимума $f(x) = -(x - 2)^2 + 5$ на целых
$x \in [0; 4]$.
\item Как представить вещественное число двоичной строкой? Как выбрать длину хромосомы по требуемой
точности?
\item В чём преимущество кода Грея? Опишите прямое и обратное преобразования. Почему неудачные запуски
в примере остановились в точке $t = 4,375$?
\item Зачем нужны мутация и элитизм? Почему в конце поиска рулетка теряет давление отбора?
\end{questions}
